% Transcription — fonds Alexandre Grothendieck, shelfmark 2, batch 3 (pages 41-44).
% Established from the facsimile archives/batches/2/batch-03.pdf (a mirror of
% https://grothendieck.umontpellier.fr/2.pdf), per the skill
% transcribe-grothendieck. The batch file carries no cover sheet and opens at
% the archivists' page 41: PDF page N is archivists' page N+40. Checked against
% the pencilled numbers bottom-left ("41", "42", "43", "44" read on the four
% leaves).
%
% Pass: Opus 5 (claude-opus-5), 2026-09-19 — first pass, unchecked against the
% pages by a human.
%
% What the batch is. Four leaves, one continuous run, the end of the folder and
% the direct continuation of page 40 (the last page of batch 2): the same ink,
% the same hand, the same notation, and the formula the first leaf opens with is
% numbered (1 ter), following the (1) and (1 bis) of page 40. Classical
% Riemann-Roch throughout — Todd class, Chern character, f_! and f_* — with no
% derived category and no perfect complex anywhere, so these pages belong with
% the campaign of pages 39-40 and not with the one of pages 32-37.
%
%   Page 41 closes the fibration example begun on page 40: the boxed (1 ter),
%     then the check that taking the degree-zero part gives back the ordinary
%     R.R. on a fibre, by computing i^*[tau_h(X).phi_X(E)] down to
%     tau_h(Z).phi_Z(E|Z); then the observation that the components of degree
%     >= 1 give Chern classes of the alternating sum, which Hirzebruch's formula
%     proper does not contain, and a footnote that for n large the Chern class
%     of f_*(E(n)) depends polynomially on n.
%   Page 42 computes phi_X(E(n)) = phi_X(E) e^{n xi} and expands
%     f_*(tau_h(X) phi_X(E(n))) as a power series in n, then opens a new
%     heading — proper morphisms, case by case — with 1) a curve mapped into a
%     variety X, K(X) = Z + Pic, xi = (n,d), Todd(X) = -k/(1-e^k) = 1 - k/2,
%     and, for Y a curve, the boxed f_!(xi) = (n delta, n delta k'/2 + f_*(d) -
%     n f_*(k/2)).
%   Page 43 draws the covering example off that formula, then does the case
%     b) Y a non-singular surface: Todd(Y) and its inverse to second order,
%     f_*(1) = delta C, and two boxed results — phi_Y(f_!(xi)) and its
%     exponential c_Y(f_!(xi)) — with the specialisation n = delta = 1.
%   Page 44 is six lines and then stops, the rest of the leaf blank: the top
%     Chern class of the structure sheaf of a curve C in the surface Y,
%     c^2(O_C) = (1/2)[C^2 - C k_Y + f_*(k_{C tilde})], and its degree
%     (1/2)(C^2 - C k_Y) + g - 1 for Y compact and the normalisation of genus g.
%
% No page is skipped: all four carry mathematics in his hand.
%
% Notation fixed once here and held, and continuous with batch 2. The Todd class
% is \tau_{h} (of a variety, \tau_{h}(X), or of a Chern class, \tau_{h}(C_{Z}(T_{Z})));
% the Chern character is \varphi, indexed by the base, \varphi_{X}, \varphi_{Y},
% \varphi_{Z}; the total Chern class is c, indexed likewise, and \widetilde{C}
% where he keeps the tilde of page 40. The direct images are f_{*} on the
% cohomological side and f_{!} on the K-theoretic one. The higher direct images
% are E^{(i)}, and E^{(0)} = f_{*}(E). The canonical class is k, lower case, k
% on X and k' on Y where both occur; on pages 43-44 he also writes k_{Y}, k_{X},
% k_{\widetilde{C}} with the base as a subscript, and both forms are kept as
% written. He underlines the divisor classes (d, k, k') on the leaf; the
% underlining is not reproduced, since \uncertain{} already sets a rule under
% what it marks and the two would be indistinguishable. The degree of f is
% \delta, the class of the image curve C, the genus g. On page 42 the class of
% the twisting sheaf is \xi and E(n) = E \otimes (that class to the n). K_{p} is
% his operator taking the component of degree p, introduced on page 41 and used
% again as K_{2} on page 44.
%
% The letters vacillate, and are kept as written. Pages 42-44 blot out a letter
% at several places, apparently while renaming source and target: page 42 heads
% the run "d'une courbe [struck] dans une variété X" and then writes K(.) = Z +
% P(.) with both arguments inked over, while every formula from there on treats
% X as the curve and Y as the target, and the boxed result carries f : X -> Y.
% Page 43 calls f_{*}(O_{X}) "un fibré vectoriel sur X" where the calculation
% puts it on Y, and page 44 calls the surface Y in the formulas and X in the
% sentence between them. None of this is repaired in the transcription; each
% occurrence carries a \note.
%
% Readings flagged and not settled. On page 41 the word after "Prenant les",
% which governs the whole reduction; the parenthesis after "R.R. usuel", read
% "(Y réduit à un point !)"; the interlinear insertion above "(L" before
% "formule de Hirzebruch"; the qualifier in parentheses after "classes de
% Chern"; the parenthesis before "de E^{(0)}"; and, in the footnote running
% along the foot of the leaf under a rule, the opening words read "pour n
% grand" and the adjective read "modifiée". The marginal note up the lower left
% edge is legible only in part — "ensuite" and "dans le cas général" — and its
% first words are lost. On page 42 the second sentence of the opening line; the
% interlinear insertion over "courbe"; the verb after "On va évidemment"; the
% word read "étant" before "précisément donné"; and the struck word replaced by
% the interlinear read "calculer". On page 43 the qualifier after "revêtement",
% the degree of the bundle f_{*}(O_{X}), the three struck lines opening case b)
% — of which only "classe canonique" and the two k_{Y} come through — the sign
% joining the two terms in the bracket of tau_h(Y)^{-1}, and the words defining
% the class C. On page 44 the target of f : C tilde -> ., inked over, and the
% preposition read "dans" before X.
\documentclass[11pt,a4paper]{article}
% Le préambule commun à toutes les transcriptions du fonds Grothendieck.
%
% Il tient en une idée : **l'incertitude se compose, elle ne se résout pas.**
% Une transcription qui lisse un mot douteux a détruit la seule chose pour
% laquelle on la faisait — un lecteur doit pouvoir savoir, à chaque endroit,
% ce qui était sur le feuillet et ce qui a été ajouté. Les six macros qui
% suivent sont donc l'appareil critique, et non de la décoration.
%
% Les mêmes commandes sont reconnues par `scripts/render.mjs`, qui produit la
% vue de lecture du site. Une macro ajoutée ici doit l'être aussi là-bas,
% faute de quoi le rendu échouera bruyamment — ce qui est voulu : un rendu
% silencieusement faux est pire qu'un rendu absent.

\ProvidesPackage{grothendieck}[2026/08/08 Transcriptions du fonds Grothendieck]

\usepackage{amsmath,amssymb,amsthm}
\usepackage{mathrsfs}
\usepackage{stmaryrd}
% Les diagrammes commutatifs sont partout dans ce fonds : c'est la langue
% même de la géométrie algébrique de Grothendieck, et une transcription qui
% les rendrait en prose ne transcrirait rien.
\usepackage{tikz-cd}
% Français par défaut — c'est la langue des feuillets — mais l'anglais est
% chargé aussi : la traduction et le résumé sont des documents anglais, et la
% typographie française (espaces avant les deux-points) y serait une faute.
% Ces documents basculent par \selectlanguage{english} en tête de corps.
\usepackage[english,french]{babel}
\usepackage{fontspec}
\usepackage{geometry}
\geometry{a4paper,margin=25mm,marginparwidth=28mm}
\usepackage{marginnote}
\usepackage{xcolor}
% ulem, not soul. `soul`'s \st and \ul are built on a character-by-character
% scan that XeLaTeX's font machinery breaks: the document fails with an
% undefined control sequence rather than degrading. ulem does the same two jobs
% and is Unicode-engine safe. `normalem` stops it hijacking \emph, which the
% transcriptions use for Grothendieck's own emphasis.
\usepackage[normalem]{ulem}
\usepackage{hyperref}
\hypersetup{colorlinks=true,linkcolor=black,urlcolor=black,pdfborder={0 0 0}}

% --- Métadonnées du lot -----------------------------------------------------
% Reprises telles quelles par le script de rendu, qui les affiche en tête de
% la vue de lecture. Elles fixent ce que le fichier prétend couvrir : sans
% elles, une transcription flotte sans point d'attache dans un fonds de seize
% mille feuillets.

\newcommand{\folder}[1]{\gdef\grothFolder{#1}}
\newcommand{\batch}[1]{\gdef\grothBatch{#1}}
\newcommand{\leaves}[2]{\gdef\grothFirst{#1}\gdef\grothLast{#2}}
\newcommand{\foldertitle}[1]{\gdef\grothFoldertitle{#1}}
\gdef\grothFolder{?}\gdef\grothBatch{?}\gdef\grothFirst{?}\gdef\grothLast{?}\gdef\grothFoldertitle{}

% --- Le résumé --------------------------------------------------------------
%
% Il ouvre la lecture modernisée, sous le titre « Résumé », et il est composé
% en petit. Ce n'est pas un ornement : c'est le seul passage du document qui
% ne soit pas des mathématiques, et un lecteur doit pouvoir le reconnaître
% avant de l'avoir lu — puis le sauter s'il connaît déjà le sujet, ou s'y
% arrêter s'il ne le connaît pas. Le filet qui le suit marque la reprise du
% corps.

\newenvironment{resume}
  {\small\setlength{\parskip}{0.45em}}
  {\par\medskip\hrule\medskip}

% --- Mots-clés ---------------------------------------------------------------
%
% En anglais, en clôture du résumé de la lecture modernisée : le vocabulaire
% moderne sous lequel chercher ce que les feuillets construisent. Le manifeste
% du site (`npm run manifest`) les extrait pour étiqueter la cote entière —
% cette ligne est la source unique des tags, il n'y a pas de fichier à côté.
\newcommand{\keywords}[1]{\par\smallskip\noindent
  {\footnotesize\textsc{Keywords} --- \textit{#1}}\par}

% --- L'appareil critique ----------------------------------------------------

% Le numéro de feuillet, en marge. C'est l'ancre qui permet de régler un
% désaccord sur une formule en regardant *un* feuillet plutôt qu'une chemise
% de six cents. Le site s'en sert aussi pour tourner les pages du fac-similé
% quand on fait défiler la transcription.
\newcommand{\leaf}[1]{%
  \marginnote{\footnotesize\color{gray!70}\textsf{#1}}%
  \label{leaf:#1}\ignorespaces}

% Illisible : on ne devine pas. Un mot inventé qui se lit comme les autres est
% le pire résultat possible.
\newcommand{\ill}{\textcolor{red!60!black}{[\,\ldots\,]}}

% Lecture douteuse : le mot est donné, et signalé comme incertain.
\newcommand{\uncertain}[1]{\uline{#1}}

% Ajout de l'éditeur — une lettre manquante, un mot sous-entendu.
\newcommand{\add}[1]{\textcolor{blue!50!black}{[#1]}}

% Biffé par Grothendieck lui-même. Ce qu'il a rayé fait partie du document :
% une rature dit souvent où la pensée a changé de direction.
\newcommand{\struck}[1]{\sout{#1}}

% Note du transcripteur, sur l'état matériel du feuillet.
\newcommand{\note}[1]{\par\smallskip{\small\color{gray!60!black}\textit{[#1]}}\par\smallskip}

% Note marginale de Grothendieck, distincte du corps du texte.
\newcommand{\marginal}[1]{\par\smallskip\noindent\hspace{1em}%
  {\small\color{gray!60!black}#1}\par\smallskip}

% --- Titre ------------------------------------------------------------------

\AtBeginDocument{%
  \begin{center}
    {\large\bfseries Fonds Alexandre Grothendieck --- cote n\textsuperscript{o}~\grothFolder}\\[2pt]
    {\itshape\grothFoldertitle}\\[4pt]
    {\small feuillets \grothFirst--\grothLast\ (lot \grothBatch)}\\[2pt]
    {\footnotesize\color{gray!60!black}Transcription établie d'après le fac-similé numérisé
     par l'Université de Montpellier.\\
     Les lectures incertaines sont soulignées, les illisibles notées [\,\ldots\,],
     les ajouts entre crochets.}
  \end{center}
  \vspace{1em}}

\endinput


\folder{2}
\batch{3}
\pages{41}{44}
\dating{[1957]}
\watermark{Édition de démonstration}
\foldertitle{Théorème de Riemann-Roch / RR [Riemann-Roch] Notes Antiques (antérieur[es à] 1958) : notes manuscrites (s.d.).}

\begin{document}

\section*{Le cas d'une fibration, et les premiers exemples}

\page{41}
\note{la page commence par la formule encadrée numérotée (1 ter), qui poursuit
les formules (1), (1 bis) et (2) de la page 40 : les deux feuillets sont d'une
même course}

\[
  \tau_{h}(Y)\sum_{i}(-1)^{i}\,\varphi_{Y}\bigl(E^{(i)}\bigr)
  \;=\; f_{*}\bigl(\tau_{h}(X)\cdot\varphi_{X}(E)\bigr)
  \tag{1 ter}
\]

Prenant les \ill{} des deux membres, ce\add{tte} formule se réduit
essentiellement à : R.R. usuel ($Y$ \uncertain{réduit} à un
\uncertain{point} !) sur le fibre. En effet, cet énoncé étant local, le membre
de gauche est égal :
\[
  \sum_{i}(-1)^{i}\,\mathrm{rg}\,E^{(i)} \;=\; \chi(Z;\,E|Z)
\]
($Z$ étant une fibre de $X \to Y$), le membre de droite est égal aussi :
\[
  K_{p}\Bigl(i^{*}\bigl[\bigl(\tau_{h}(X)\cdot\varphi_{X}(E)\bigr)^{(p)}\bigr]\Bigr),
\]
où $i : Z \to X$ \ill{} l'application identique, et où $p$ \ill{} la dimension
de la fibre, et où $^{(p)}$ \ill{} \ill{} \ill{} de degré $p$. On a \ill{} :
\begin{align*}
  i^{*}[\;\;] &= i^{*}\bigl(\tau_{h}(X)\bigr)\,\struck{\ill{}}\cdot
                 i^{*}\bigl(\varphi_{X}(E)\bigr) \\
              &= \tau_{h}\bigl(C_{Z}(i^{*}(T_{X}))\bigr)\,
                 \varphi_{Z}\bigl(i^{*}(E)\bigr) \\
              &= \tau_{h}\bigl(C_{Z}(T_{Z})\bigr)\,\varphi_{Z}(E|Z)
               \;=\; \tau_{h}(Z)\,\varphi_{Z}(E|Z) .
\end{align*}

Ce\add{ci} redonne bien \add{(}\ill{}\add{)} la formule de Hirzebruch pour le
fibre. --- Mais, prenant des \struck{\ill{}} composantes de degré
$\geqslant 1$ des deux membres, on obtient des classes de Chern \add{(}\ill{}\add{)}
de $\sum_{i}(-1)^{i}E^{(i)}$ \add{(}\ill{}\add{)} de $E^{(0)}$, ce qui n'est
pas contenu dans la formule de Hirzebruch proprement dite. Ce\add{ci}
\uncertain{traduit} en particulier le résultat suivant (qui n'est pas non plus
trivial) : \uncertain{pour $n$ grand}, la classe de Chern \uncertain{modifiée}
de $E(n)^{\circ} = f_{*}(E(n))$ dépend polynômialement \ill{}
\note{les deux dernières lignes courent au bas du feuillet sous une longue
barre, terminée par une flèche vers la droite : c'est une note ajoutée au
résultat annoncé, et non la suite du texte}

\marginal{\ill{} ensuite \ill{} dans le cas général}

\page{42}
de $n$. \ill{}
\note{la ligne qui ouvre le feuillet achève la note portée au bas de la page
41 : « dépend polynômialement » / « de $n$ ». Ce qui la suit, d'une encre plus
pâle, ne se laisse pas lire}

\[
  \varphi_{X}\bigl(E(n)\bigr) \;=\; \varphi_{X}(E)\,\struck{\ill{}}\,e^{n\xi} .
\]
d'où
\[
  \tau_{h}(X)\,\varphi_{X}\bigl(E(n)\bigr)
  \;=\; \tau_{h}(X)\,\varphi_{X}(E)\,e^{n\xi}
  \;=\; \sum_{k} n^{k}\left[\frac{\xi^{k}}{k!}\,\tau_{h}(X)\,\varphi_{X}(E)\right]
\]
et
\[
  f_{*}\bigl(\tau_{h}(X)\,\varphi_{X}(E(n))\bigr)
  \;=\; \sum_{k} n^{k}\, f_{*}\!\left(\frac{\xi^{k}}{k!}\,\tau_{h}(X)\,\varphi_{X}(E)\right)
\]

\subsection*{Cas de morphismes propres particuliers}
\note{le titre est de lui, séparé de ce qui précède par un trait ondulé}

1) Morphismes d'une courbe \struck{\ill{}} dans une variété $X$.
\note{une addition interlinéaire d'une encre plus claire court au-dessus de
« courbe » et rejoint le mot biffé : elle ne se laisse pas lire. À partir
d'ici et jusqu'au bas du feuillet, $X$ désigne la courbe et $Y$ le but ; la
lettre est effacée sous un pâté à chacune de ses premières occurrences, et la
formule encadrée porte $f : X \to Y$}

On va évidemment \uncertain{supposer} $f$ un \uncertain{morphisme}, $X$ et $Y$
irréductibles.

\struck{\ill{}}
\[
  K(\ill{}) \;\simeq\; \mathbf{Z} + P(\ill{}),
\]
un isomorphisme \uncertain{étant} précisément donné par le morphisme
\struck{\ill{}} $\widetilde{C}$ de Chern. On doit \ill{} \struck{\ill{}}
\uncertain{calculer}, où $\xi = (n,d)$ \ill{} $f_{!}(\xi) \in K(Y)$ \ill{}
\uncertain{fonction} de $n$, $d$. Notons \ill{}
\[
  \varphi_{X}(\xi) = n + d,
  \qquad \widetilde{C}(T_{X}) = (1, -k),
\]
d'où
\[
  \tau_{h}(X) \;=\; \frac{-k}{1 - e^{k}}
  \;=\; 1 - \tfrac{1}{2}k,
\]
$k$ étant la classe canonique. D\ill{} \struck{\ill{}}.
\[
  \varphi_{X}(\xi)\,\tau_{h}(X)
  \;=\; \bigl(1 - \tfrac{1}{2}k\bigr)(n + d)
  \;=\; \struck{\ill{}}\;\; n + \bigl(d - n\tfrac{1}{2}k\bigr)
\]

a) $Y$ est une courbe. Soit $k'$ son diviseur canonique, d'où
$\tau_{h}(Y) = 1 - \tfrac{1}{2}k'$. Soit $\delta$ le degré de $f$.
Alors on a $f_{*}(1) = \delta$, donc
\[
  f_{*}\bigl(\varphi_{X}(\xi)\,\tau_{h}(X)\bigr)
  \;=\; n\delta + \bigl(f_{*}(d) - n f_{*}(\tfrac{1}{2}k)\bigr)
\]
d'où
\begin{align*}
  \varphi_{X}\bigl(f_{!}(\xi)\bigr)
    &= \tau_{h}(Y)^{-1} f_{*}(\cdot\;\ill{})
     = \bigl(1 + \tfrac{1}{2}k'\bigr)
       \Bigl(n\delta + \bigl(f_{*}(d) - n f_{*}(\tfrac{1}{2}k)\bigr)\Bigr) \\
    &= n\delta + \Bigl[n\delta\tfrac{1}{2}k' + f_{*}(d)
       - n f_{*}(\tfrac{1}{2}k)\Bigr]
\end{align*}
\note{le membre de gauche est écrit $\varphi_{X}$ sur le feuillet, là où le
calcul demande le caractère de Chern sur $Y$ ; la lettre est reproduite telle
quelle}
\ill{}
\[
  f_{!}(\xi) \;=\; \Bigl(n\delta,\;\;
    n\delta\tfrac{1}{2}k_{y} + f_{*}(d)
    - n f_{*}\bigl(\tfrac{1}{2}k_{x}\bigr)\Bigr)
  \qquad (f : X \to Y)
\]
\note{dans l'encadré il distingue les deux classes canoniques en portant la
base en indice : $k_{y}$ pour le $k'$ des lignes
précédentes, $k_{x}$ pour le $k$ de la courbe}

\page{43}
p. ex. si \struck{\ill{}} est un revêtement \ill{} \add{(}\ill{}\add{)} de
degré $\delta$ de $Y$ \uncertain{par} $f$, alors $f_{*}(\mathcal{O}_{X})$ est
un \struck{\ill{}} fibré vectoriel sur $X$ correspondant à un fibré
\struck{\ill{}} de degré \ill{} $-\,g\,\delta$, et donc la classe de Chern est
\[
  \tfrac{1}{2}\,\struck{\ill{}}\bigl(\delta k_{Y} - f_{*}(k_{X})\bigr) .
\]
\note{le feuillet place $f_{*}(\mathcal{O}_{X})$ « sur $X$ », là où le calcul
le demande sur $Y$ ; la lettre est reproduite telle quelle}

b) $Y$ est une surface non singulière. Soient $k_{Y}$ et \ill{} les deux
\ill{} de Chern des fibrés \struck{tangents} \ill{}, dont $k_{Y}^{(1)}$ est la
classe canonique, $k_{Y}^{2}$ \struck{\ill{}}.
\note{ces trois lignes sont barrées d'un long trait horizontal et portent deux
additions interlinéaires : seuls les mots donnés ici se laissent lire}

\[
  c(T_{Y}) \;=\; 1 - k_{Y}^{1} + \ill{}
\]
d'où
\[
  \tau_{h}(Y) \;=\; 1 - \tfrac{1}{2}k_{Y}^{1} + \struck{\ill{}}\;\ill{},
\]
et \ill{}
\[
  \tau_{h}(Y)^{-1} \;=\; 1 + \tfrac{1}{2}k_{Y}^{1} + \struck{\ill{}}
  + \tfrac{1}{4}k_{Y}^{1\,2}
  \;=\; 1 + \tfrac{1}{2}k_{Y}^{1}
  + \tfrac{1}{12}\bigl[\,k_{Y}^{2}\;\ill{}\;k_{Y}^{1\,2}\,\bigr]
\]
\note{le signe qui joint les deux termes du crochet est recouvert et ne se
laisse pas trancher}

\[
  f_{*}(1) \;=\; \delta f(X) \;=\; \delta\,\struck{\ill{}}\;C
\]
$C$ \ill{} de $X$ \ill{} définie comme classe de diviseurs, et $\delta$ le
degré de $X \to C$.

$f_{*}(k_{X})$, $f_{*}(d)$ : classes de $0$-cycles.

\[
  \varphi_{Y}\bigl(f_{!}(\xi)\bigr)
  \;=\; \Bigl[\,1 + \tfrac{1}{2}k_{Y} + \struck{\ill{}}\,\Bigr]
        \Bigl[\,\struck{\ill{}}\;\; n\delta C
          + \bigl(f_{*}(d) - \tfrac{1}{2}n\,f_{*}(k_{x})\bigr)\Bigr]
\]
\[
  \varphi_{Y}\bigl(f_{!}(\xi)\bigr) \;=\; n\delta C
  + \Bigl[\,f_{*}(d) - \tfrac{1}{2}n\,f_{*}(k_{x})
    + \tfrac{1}{2}n\delta\,(C\cdot k_{Y})\,\Bigr]
\]
\begin{align*}
  c_{Y}\bigl(f_{!}(\xi)\bigr)
    &= \exp\bigl(n\delta C - [\;\;]\bigr) \\
    &= 1 + n\delta C - [\;\;] + \tfrac{1}{2}n^{2}\delta^{2}C^{2}
\end{align*}
\note{le crochet laissé vide renvoie au crochet de la formule encadrée qui
précède}
d'où
\[
  c_{Y}\bigl(f_{!}(\xi)\bigr) \;=\; 1 + n\delta C
  + \Bigl[\,\tfrac{1}{2}n^{2}\delta^{2}C^{2}
    - \tfrac{1}{2}n\delta\,C k_{Y} - f_{*}(d)
    + \tfrac{1}{2}n\,f_{*}(k_{X})\,\Bigr]
\]
\ill{}, $n = \delta = 1$ \struck{\ill{}} on trouve
\[
  c_{Y}\bigl(f_{!}(\xi)\bigr) \;=\; 1 + C
  + \Bigl[\,\tfrac{1}{2}\bigl(C^{2} - C k_{Y}\bigr) - f_{*}(d)
    + \tfrac{1}{2}f_{*}(k_{X})\,\Bigr]
\]
\note{sur le feuillet les trois résultats de ce paragraphe sont encadrés}

\page{44}
En particulier, on trouve la dernière classe de Chern d'une courbe $C$ dans la
surface $Y$ :
\[
  c^{2}(\mathcal{O}_{C}) \;=\; \tfrac{1}{2}\bigl[\,C^{2} - C k_{Y}
    + f_{*}(k_{\widetilde{C}})\,\bigr]
\]
\note{l'indice du dernier $k$ est recouvert d'encre ; il est lu sur la ligne
suivante, qui nomme $k_{\widetilde{C}}$}
où $f : \widetilde{C} \to \ill{}$ est l'homomorphisme canonique
\uncertain{dans} $X$ de la \struck{\ill{}} normalisée $\widetilde{C}$ de $C$,
et $k_{\widetilde{C}}$ la classe canonique sur $\widetilde{C}$. Si $Y$ est
\uncertain{compacte}, et $\widetilde{C}$ \ill{} genre $g$, on trouve, en
prenant le $K_{2}$ des deux membres :
\[
  K_{2}\bigl(c^{2}(\mathcal{O}_{C})\bigr)
  \;=\; \tfrac{1}{2}\bigl[\,C^{2} - C k_{Y} + 2g - 2\,\bigr]
  \;=\; \tfrac{1}{2}\bigl(C^{2} - C k_{Y}\bigr) + g - 1
\]
\note{la surface est appelée $Y$ dans les formules et $X$ dans la phrase qui
les commente, sans que rien sur le feuillet distingue les deux ; l'opérateur
$K_{2}$ est celui de la page 41, où $K_{p}$ prend la composante de degré $p$.
Le feuillet s'arrête ici et le reste en est blanc : c'est la fin du dossier}

\end{document}
